\section*{Chapter \ref{ch:dv:monte-carlo-pricers-in-practice} --- Monte Carlo Pricers in Practice}

\begin{solution}{exo:dv:monte-carlo-pricers-in-practice:1}
The Bermudan's price is the supremum, over stopping rules, of the expected discounted exercise value. The fitted policy is one stopping rule, so its
expected value is at most the price. On paths independent of the fit, the policy is fixed before the paths are drawn, and the average is an unbiased
estimate of that policy's value, which is a lower bound.
\end{solution}

\begin{solution}{exo:dv:monte-carlo-pricers-in-practice:2}
$S_T=S_0\exp\bigl((r-\tfrac12\sigma^2)T+\sigma\sqrt T Z\bigr)$, so $\partial S_T/\partial S_0=S_T/S_0$. The payoff $(S_T-K)^+$ has derivative
$\mathbf 1\{S_T>K\}$ except at $S_T=K$, which has probability zero. The chain rule then gives $e^{-rT}\mathbf 1\{S_T>K\}S_T/S_0$ on each path, and its
average is the delta.
\end{solution}

\begin{solution}{exo:dv:monte-carlo-pricers-in-practice:3}
The digital's payoff is flat on both sides of the strike, so its derivative is zero on almost every path. The delta comes entirely from probability mass
crossing the strike, which the pathwise method, differentiating inside the expectation, cannot see; the interchange of derivative and expectation fails
for a discontinuous payoff. The likelihood ratio differentiates the density of $S_T$, which is smooth, and keeps the payoff as it is.
\end{solution}

\begin{solution}{exo:dv:monte-carlo-pricers-in-practice:4}
$x=\ln S_T$ is normal with mean $\mu=\ln S_0+(r-\tfrac12\sigma^2)T$ and standard deviation $\sigma\sqrt T$. Its log-density is
$-\ln(\sigma\sqrt T)-\tfrac12Z^2+\text{const}$ with $Z=(x-\mu)/(\sigma\sqrt T)$. Differentiating in $\sigma$ with $x$ fixed:
$\partial Z/\partial\sigma=-Z/\sigma-(\partial\mu/\partial\sigma)/(\sigma\sqrt T)=-Z/\sigma+\sqrt T$, so the score is
$-1/\sigma-Z(-Z/\sigma+\sqrt T)=(Z^2-1)/\sigma-Z\sqrt T$.
\end{solution}

\begin{solution}{exo:dv:monte-carlo-pricers-in-practice:5}
Given the outer path, the inner estimates are the true continuation values plus noise $\varepsilon_j$ of mean zero. The bound averages
$\max_j(a_j+\varepsilon_j)$, where $a_j$ is what the bound would be with exact values. The maximum is convex, so by Jensen's inequality
$\E[\max_j(a_j+\varepsilon_j)]\ge\max_j a_j$. Noise can only raise the bound, and more inner paths shrink the excess, as the put's
6.134, 6.054 and 6.037 show.
\end{solution}

\begin{solution}{exo:dv:monte-carlo-pricers-in-practice:6}
Pseudo-random coordinates are independent and identically distributed. Which normal drives which piece of the path does not change the estimator's
distribution, so the bridge changes nothing. Low-discrepancy coordinates are not all equally good: the first ones are the best spread, and in high
dimensions the later ones are poorly spread in pairs. The bridge puts most of the path's variance into the first coordinates, lowering the effective
dimension of the integrand.
\end{solution}

\begin{solution}{exo:dv:monte-carlo-pricers-in-practice:7}
With 2\,000 training paths, averaged over 40 training sets, fitting and pricing on the same paths gives 18.74; the same policies applied to 100\,000
independent paths give 18.50. The difference is $+0.25$ (standard error 0.06), and the in-sample price is above even the \omterm{def:dv:monte-carlo-pricers-in-practice:dual}{dual upper bound} 18.713. The
regression has seen each path's future, so on those paths it exercises with foresight: the bias is upwards. With 200\,000 training paths and nine
regressors, the bias falls below the noise.
\end{solution}

\begin{solution}{exo:dv:monte-carlo-pricers-in-practice:8}
The standard error measures only the noise of the lower-bound estimate. The price is biased low by the policy's suboptimality, and the standard error
does not see that bias: with the small basis the standard error is 0.039 but the bounds are 0.668 apart. Two cents needs the dual bound as well (and a
check of the time discretisation where the model has one).
\end{solution}

\begin{solution}{pb:dv:monte-carlo-pricers-in-practice:1}
\textbf{1.} 18.649, standard error 0.039.
\textbf{2.} Pricing on the fitting paths is biased upwards (exercise 7: $+0.25$ with 2\,000 training paths); independent paths make the price a lower bound.
\textbf{3.} 18.137.
\textbf{4.} It sees only the leader. When the second price is close to the first, the option on the maximum is worth more alive, and the small basis cannot
tell those states apart; it exercises too soon there, and too late where the leader is far ahead.
\textbf{5.} 15.57 (standard error 0.07): a lower bound for the Bermudan. The \omterm{def:dv:american-options-and-early-exercise:premium}{early-exercise premium}, about 3.08, comes from the 10\% dividend yield, which the
holder gives up by waiting.
\textbf{6.} 18.713, standard error 0.019.
\textbf{7.} 0.064, 0.3\% of the price.
\textbf{8.} 18.805; the gap is 0.668, ten times wider.
\textbf{9.} 6.134, 6.054 and 6.037 against the tree's 6.0326: the excess over the tree falls about fivefold each time the inner paths are multiplied by four.
\textbf{10.} Each outer node needs its own inner simulation: $1\,000\times1\,000$ inner paths started at each of nine dates, running out to the last date,
are 45 million path steps, against 1.8 million for the lower bound's 200\,000 paths, 25 times more.
\textbf{11.} Pathwise, holding the fitted policy fixed: at the optimal policy the value's sensitivity to the boundary is zero to first order, so fixing it costs
little. Or bump with common random numbers and the same policy.
\textbf{12.} The difference of two independent estimates has the variance of both divided by the square of the bump, huge for a small bump. Refitting the
policy on new paths adds a jump of its own.
\textbf{13.} All sensitivities (three spots, three volatilities, correlations, rates) of every trade for a cost of a few pricings, independent of their number.
\textbf{14.} The path is a 27-dimensional point (three assets at nine dates); a bridge per asset and a principal-component ordering across assets lower the
effective dimension. The exercise decision is a discontinuity, which limits the gain.
\textbf{15.} None in time: lognormal prices are simulated exactly at the exercise dates. Local or stochastic volatility would need small steps between the
dates and bring a bias.
\textbf{16.} The desk is short, and the holder may exercise better than the desk's policy, so the liability is up to the upper bound. Quote around the bounds'
midpoint and book the lower bound with a reserve for the gap.
\textbf{17.} The gap plus about two standard errors of the bounds, $0.064+0.086\approx0.15$ per 100 of notional; with the small basis it would be about 0.8.
\textbf{18.} When the gap is within the statistical noise of the bounds: 0.064 against two combined standard errors of 0.086, so the rich basis has reached it.
\textbf{19.} Rich basis: 18.649 against 18.713, a gap of 0.064 (0.3\% of the price). Small basis: 18.137 against 18.805, a gap of 0.668, ten times wider.
\textbf{20.} The most the price can be wrong because of the exercise policy, plus the upward bias of the inner simulation.
\end{solution}

\begin{solution}{iq:dv:monte-carlo-pricers-in-practice:1}
Simulate paths forward. Backwards from the last exercise date, regress the discounted realised cash flows on basis functions of the state, using
in-the-money paths, and exercise where the exercise value beats the fitted continuation; carry the realised cash flows back. Then apply the policy to new
paths: that price is a lower bound.
\iqlookfor{in-the-money regression, realised cash flows, independent pricing paths, lower bound.}
\end{solution}

\begin{solution}{iq:dv:monte-carlo-pricers-in-practice:2}
By duality: the price is at most $\E[\max_j(h_j-M_j)]$ for any martingale $M$ with $M_0=0$, with equality for the right one. Andersen--Broadie build $M$ from a
lower-bound policy with nested simulation of its continuation values. The gap between the bounds measures the policy's quality.
\iqlookfor{the dual formula, nested simulation, the gap.}
\end{solution}

\begin{solution}{iq:dv:monte-carlo-pricers-in-practice:3}
Pathwise differentiates the payoff along each path: low variance, but it needs a payoff continuous in the parameter; it fails for digitals and barriers.
Likelihood ratio differentiates the density: any payoff, but higher variance (6.5 times for a call's delta in the chapter), and it needs a known density,
awkward after many steps. Smoothing or conditioning mixes the two.
\iqlookfor{continuity versus density, and the variance trade-off.}
\end{solution}

\begin{solution}{iq:dv:monte-carlo-pricers-in-practice:4}
Andersen's \omterm{def:dv:monte-carlo-pricers-in-practice:qe}{quadratic-exponential scheme} for the variance, with the integrated variance from both ends of the step for the log-price, and a martingale
correction if needed. Euler makes the variance negative when the Feller condition fails, and truncation leaves a bias that decays slowly (0.86 at four
steps a year in the chapter, against QE within the noise).
\iqlookfor{QE, and why Euler's truncation is biased.}
\end{solution}

\begin{solution}{iq:dv:monte-carlo-pricers-in-practice:5}
Fix the seeds and use the same random numbers for the base and bumped runs; hold the exercise policy fixed; replace bumps by pathwise, likelihood-ratio or
adjoint \omterm{def:dv:greeks-and-the-hedging-pnl:greeks}{Greeks}; smooth discontinuous payoffs; add control variates; check that the noise, not the market, moves the numbers.
\iqlookfor{common random numbers and direct \omterm{def:dv:greeks-and-the-hedging-pnl:greeks}{Greeks}.}
\end{solution}

\begin{solution}{iq:dv:monte-carlo-pricers-in-practice:6}
Run the paths forward and store the states; then run backwards along each path, propagating the derivative of the payoff with respect to each state (the
adjoint) and accumulating the sensitivities to every input. All \omterm{def:dv:greeks-and-the-hedging-pnl:greeks}{Greeks} cost a small multiple of one pricing, whatever their number; the price is memory for
the stored paths (or recomputation by checkpointing), and a pathwise method's need for continuous payoffs.
\iqlookfor{the reverse pass, cost independent of the number of inputs, memory.}
\end{solution}
